\clearpage
\section{Full forward operational robustness and converse classification}
\label{sec:copying-operational-forward-robustness}
This section completes the allocation and assignment residuals, the general
deep beta/copying cases, and the classification of arbitrary converse overlaps
identified in the preceding milestone. It also assembles universal forward
robustness for both the micro-beta and reference-creation templates. The proof
uses the original copying machine, its complete operation alphabet, the
original stability and copyability premises, and the binding-free coordinate
specimen system. No operational matching, determinism, local-diagram or
residual-admissibility premise is left unproved in the forward theorem.
The Proposition~12.2 terms and their Figure~16 translation are unchanged.

\begin{theorem}[Universal forward operational robustness]
Let $\mathcal T$ be either the micro-beta or reference-creation pre-template,
and let $Q,P,U$ be cost orders. Its original binding-free specimen system
satisfies $\robust{Q}{P}{U}$. In the substantive matching case, a supplied
left successor is matched by one actual right rewrite and an admissible
pointed quasi-specimen, using zero additional left steps. The quantitative
comparison is $Q(1+0,1)$, which follows from reflexivity.
\end{theorem}
The proof constructs the local diagrams for all nine original local rules,
lifts them to their actual whole-graph matches, and proves binding freedom
of the resulting common contexts. Complete local-occurrence extraction then
covers contextual renaming, plugging and expansion. Finally, reflection of
an arbitrary representation-class step recovers the finite rewrite; a search
step cannot start in rewrite mode. The resulting theorem is
\texttt{OperationalTemplates.template\_robust}; its proof and source link
appear below.

The converse classification is also exhaustive, but its output depends on
the rule. Arithmetic, allocation and beta give retained complete footprints.
Read, write and equality give intrinsic matches in the original common
context. Copying gives a retained complete payload and an original selected
sharing resolution. Every branch identifies the actual supplied successor,
including the resumed token. It does not assume that every unused branch of
an arbitrarily large right-hand sharing tree remains outside template
fillings. A new minimal-spine theorem constructs a smaller complete copying
source while leaving those branches at its boundary.

At this milestone, universal converse robustness was a subsequent obligation
(now discharged in Section~\ref{sec:copying-operational-converse-robustness}):
construct all opposite rewrites and admissible residual specimens from the
reflected selectors, including atom comparisons with two selected inputs.
Classification and successor coherence alone do not establish those
existence and admissibility claims. The other eight template families,
semantic weakening and Figure~22 are also unfinished. Thus the new robustness
theorem is a substantive dependency of Proposition~12.2, not a claim that the
final example has already been proved.

\Needspace{12\baselineskip}
\subsection{Guarded and thunk outputs in allocation, assignment and beta}
\termdef{binding-or-thunk-outputs}
A graph has \emph{binding-or-thunk outputs} when each external output is either
a thunk wire or is produced by an edge with a sourced binding ancestor in
that graph. The source witness belongs to the guarding edge; it need not be
an internal vertex. This condition extends the earlier guarded-output
criterion exactly where stable lambda values expose their thunk interface.
It neither deletes those outputs nor changes the stability premise.

For a stable graph, every output is an atom or a thunk. If an atom occurs in
an edge's target list, the target typing equation puts that atom in the
primitive's output signature. Inspection of the original stable primitive
signatures leaves only the instance edge: passive constants have no outputs,
and lambda has one thunk output. That instance consumes the stable graph's
sole term input and guards its own atom output. A thunk output takes the
other alternative directly. These arguments prove the output condition for
every original stable graph, including its arbitrary carrier representation.

\lean{Copying/BindingStableOutputs}{11}{14}{BindingOrThunkOutputs}
\lean{Copying/BindingStableOutputs}{17}{28}{Stable.atom_target_instance}
\lean{Copying/BindingStableOutputs}{30}{41}{Stable.bindingOrThunkOutputs}

\paragraph{The original allocation and assignment interfaces.}
Serial composition preserves a guard in its right component. Map its sourced
edge and ancestry through the actual right occurrence; the right component's
external outputs are the serial graph's original external outputs. Thunk
types are unchanged. Allocation's source is the reference head followed by
its stable argument, so this proves its output condition immediately.

Assignment has two kinds of output. Its first output is the old store content.
The original destination instance is a sourced binding edge, and the proved
instance-to-store trace supplies its ancestry to that output. Each remaining
output belongs to the unchanged stable new value and uses the atom-or-thunk
argument above. The proof respects the original order of these outputs and
the arbitrary atom-tree presentation. For a retained primitive footprint,
transporting these conditions to the common skeleton changes neither wire
types nor incidences; primitive label reflection preserves the guarding edge.

\lean{Copying/BindingStableOutputs}{44}{60}{BindingOrThunkOutputs.serialRight}
\lean{Copying/BindingStableOutputs}{62}{64}{Allocation.before_bindingOrThunk}
\lean{Copying/BindingStableOutputs}{68}{96}{Assignment.before_bindingOrThunk}
\lean{Copying/BindingStableOutputs}{100}{115}{RetainedFootprint.skeleton_bindingOrThunk}

\paragraph{Why a thunk cannot provide an unguarded route to a hole.}
The auxiliary condition \texttt{HoleTermInputs} says that every shallow hole
has exactly the input interface consisting of one term wire. It is derived,
rather than assumed, for the operational common context when its left
expansion is hole-free. At an original hole, a retained corolla would leave
that same hole in the expansion, contradicting hole freedom. A filled
corolla belongs to an original operational pair, whose input is the required
single term wire. This uses the literal input interface of micro-beta and
reference creation and is unaffected by independent endpoint isomorphisms.

Now suppose an edge consumes a thunk wire. Its source typing equation rules
out every primitive signature except a possible hole. The derived hole-input
condition rules out that exception. The edge must therefore be a box. In
particular it is binding and is not a hole; its consumed thunk supplies a
source witness. This argument uses only actual signatures and incidences,
not an extra restriction on allowed operational configurations.

\lean{Copying/BindingThunkReplacement}{12}{13}{HoleTermInputs}
\lean{Copying/BindingThunkReplacement}{17}{32}{OperationalTemplates.holeTermInputs}
\lean{Copying/BindingThunkReplacement}{36}{56}{thunk_source_binding}

\paragraph{Admissibility of the completed replacement.}
Consider a complete occurrence of a source with binding-or-thunk outputs in
a binding-free host whose holes have term inputs. An adjacency leaving the
inserted replacement and entering a surviving edge must cross one original
source output. If that output has a guard inside the old source, prefixing
a retained path to a hole by that guard gives a forbidden old path. If it is
a thunk, the surviving consumer itself is a sourced binding box. It cannot
be the final hole; its positive ancestry to that hole again contradicts
original binding freedom.

Induction backwards along an ancestry ending at a surviving hole therefore
shows that every edge of the ancestry survives from the host. When the
replacement is hole-free, every possible shallow hole is such a surviving
hole. Reflect the positive forbidden path and its sourced binding edge to
the old host to obtain a contradiction. Focus freedom and the recursively
checked box conditions finish the binding-free theorem: surviving boxes
inherit their old proofs, and boxes in the closed replacement contain no
holes. The frame theorem applies this result to the genuine complete
occurrence of the matched skeleton. It proves admissibility of the filled
graft; it does not assume that the intermediate cut is binding-free.

\lean{Copying/BindingThunkReplacement}{66}{87}{CompleteMatchCut.thunk_exit_forbidden}
\lean{Copying/BindingThunkReplacement}{89}{102}{CompleteMatchCut.thunk_ancestor_survivor}
\lean{Copying/BindingThunkReplacement}{104}{128}{CompleteMatchCut.thunk_shallow_binding}
\lean{Copying/BindingThunkReplacement}{132}{148}{CompleteMatchCut.bindingFree_replace_thunk}
\par\penalty0
\lean{Copying/BindingThunkReplacement}{152}{162}{RetainedFootprint.frame_plug_bindingFree_thunk}
\par\penalty0

\paragraph{Actual whole-graph steps and their quantitative residuals.}
For a primitive retained source, the previously proved right reconstruction
identifies the paired source with the same local rule source. Apply the
unchanged rule on both sides and use the reflexive filling witness of its
actual reduct as the local residual. The new graft theorem proves its binding
freedom. The pointed residual theorem places the two tokens at the actual
corresponding successor coordinates, and rewrite preservation supplies
hole-free endpoints. This constructs both whole-graph steps and their
binding-free quasi-specimen.

For an independently supplied successor, proved machine determinism
identifies it with the constructed successor. The quantitative conclusion
uses zero additional steps on that side and one matching step on the other;
reflexivity of the cost order supplies the comparison. Exchanging the same
residual proves the corresponding converse case without imposing symmetry
on the operational relation.

\lean{Copying/PrimitiveThunkTransfer}{23}{44}{RetainedFootprint.primitive_thunk_diamond}
\lean{Copying/PrimitiveThunkTransfer}{46}{61}{RetainedFootprint.primitive_thunk_robust_case}
\lean{Copying/PrimitiveThunkTransfer}{63}{79}{RetainedFootprint.primitive_thunk_converse_robust_case}

The public allocation and assignment theorems discharge all the auxiliary
conditions for the original operational templates. Footprint avoidance comes
from the forward classification at a non-entering common-context position;
the guard comes from the original stable operand; term-input holes follow
from the closed expansion; and the original rule's safety theorem supplies
its closed, focus-free reduct. Thus no caller supplies a residual certificate,
a paired rewrite, a thunk restriction or a graft-admissibility assumption.
The converse wrappers quantify over arbitrary right successors of this same
constructed paired footprint. They do not by themselves extract a retained
footprint from every arbitrary converse overlap.

\lean{Copying/AllocationAssignmentResiduals}{16}{31}{OperationalTemplates.ForwardResidual.allocation}
\lean{Copying/AllocationAssignmentResiduals}{34}{49}{OperationalTemplates.ForwardResidual.allocation_converse}
\lean{Copying/AllocationAssignmentResiduals}{52}{67}{OperationalTemplates.ForwardResidual.assignment}
\lean{Copying/AllocationAssignmentResiduals}{70}{85}{OperationalTemplates.ForwardResidual.assignment_converse}

\paragraph{Guards survive recursive changes inside boxes.}
A retained outer edge preserves whether its label is binding even when a
boxed body changes. At depth zero, every label is primitive and the retained
primitive equality suffices. At positive depth, a primitive label can only
match the same primitive; a box cannot become primitive, because retention
reflects primitive labels in both directions. The remaining box-to-box case
is binding on both sides. Full label matching from the local source to its
actual contribution, followed by this retained-label equivalence, transports
every guarding edge to the common skeleton. The ancestry and source/output
incidences are already identical. This removes the earlier primitive-only
restriction from guarded-output transport.

\lean{Copying/RetainedBindingGuards}{11}{29}{Search.Expansion.Retained.bindingLabel_iff}
\lean{Copying/RetainedBindingGuards}{33}{43}{RetainedFootprint.skeleton_bindingGuarded_full}

\paragraph{The copying source's guards.}
Guarded outputs transport through recursive graph isomorphisms and through
the right component of serial composition by mapping the guarding edges,
source witnesses and ancestry. A positive-input contraction tree supplies
its own guard: follow the original trace from the selected input to the
unique root. Its first edge has a source, so it cannot be a weakening and
must be a contraction. It is therefore binding. The entire original trace
connects this sourced contraction to the root output, independently of
unused branches or the tree's presentation.

If each output-producing edge of a payload also consumes the payload input,
append the actual serial adjacency from the tree root to that edge. This
proves the copying-source guard for the primitive payload cases. A lambda
payload instead has a sourced box guarding all its free environment outputs.
Induction on copyability now covers instances, constants, application, every
unary and binary operation, arbitrary lambda bodies, and recursive renaming.
Active operations are included: copyability is not replaced by stability.

\lean{Copying/CopySourceBindingGuards}{10}{18}{BindingGuardedOutputs.iso}
\lean{Copying/CopySourceBindingGuards}{20}{36}{BindingGuardedOutputs.serialRight}
\lean{Copying/CopySourceBindingGuards}{40}{59}{CopyRule.tree_bindingGuarded}
\lean{Copying/CopySourceBindingGuards}{63}{87}{CopyRule.before_flat_bindingGuarded}
\par\penalty0
\lean{Copying/CopySourceBindingGuards}{90}{96}{lambdaPayload_bindingGuarded}
\par\penalty0
\lean{Copying/CopySourceBindingGuards}{100}{125}{CopyRule.before_bindingGuarded}

\paragraph{The beta source's guards.}
A guarded output is, in particular, binding-or-thunk. This more general
condition composes over the original ordered tensor: split an output into
its actual left or right component, preserve its wire label, and map that
component's guard through its complete ordered occurrence. In a beta source,
the lambda's box guards the closure's free outputs; the stable argument
supplies guards or thunk types for its own outputs. Tensor these two results
and follow them by the original application head in serial composition.
This establishes the condition for the complete beta source without any
restriction on the lambda body. The full-label skeleton theorem transports
it through recursive box filling using the retained binding-label equivalence;
no primitive-footprint hypothesis is needed.

\lean{Copying/BetaSourceBindingGuards}{10}{12}{BindingGuardedOutputs.bindingOrThunk}
\lean{Copying/BetaSourceBindingGuards}{15}{32}{BindingOrThunkOutputs.tensor}
\lean{Copying/BetaSourceBindingGuards}{36}{41}{Application.before_bindingOrThunk}
\lean{Copying/BetaSourceBindingGuards}{46}{57}{RetainedFootprint.skeleton_bindingOrThunk_full}

\Needspace{12\baselineskip}
\subsection{Finite selected contraction spines}
\termdef{minimal-selected-spine}
A \emph{minimal selected spine} retains precisely the forward contraction
path from one selected tree input to the root. Side branches become boundary
attachments. Here ``minimal'' refers to that selected path; it is not a
claim that the original graph or machine rule is replaced by a global
normal form. The construction is uniform for term and atom sharing and
preserves each contraction's ordered input slots.

\paragraph{An explicit finite graph.}
For $n+1$ contractions, use $n$ internal vertices, $n+2$ external inputs and
one output. The first external input is selected. Edge $e$ consumes the
selected input when $e=0$, and otherwise consumes the previous internal path
vertex. Its other input is external input $e+1$. Its target is internal
vertex $e$ except at the last edge, which produces the sole external output.
A Boolean orientation records whether the selected branch occupies the
first or second ordered source slot. This defines the vertex carrier,
entry, side, output and ordered-source functions explicitly.

\lean{Copying/ContractionSpine}{12}{12}{ContractionSpine.Vertex}
\lean{Copying/ContractionSpine}{14}{15}{ContractionSpine.entry}
\lean{Copying/ContractionSpine}{16}{16}{ContractionSpine.side}
\lean{Copying/ContractionSpine}{17}{18}{ContractionSpine.out}
\par\penalty0
\lean{Copying/ContractionSpine}{20}{21}{ContractionSpine.sources}
\par\penalty0

Membership in a source row is independent of orientation: it means equality
to the entry or side vertex. Entries are distinct from all side vertices,
and the entry, side and output maps are separately injective. These facts
prove uniqueness of the consuming and producing edges. Coverage is checked
by port cases. An internal vertex is consumed by the next edge and produced
by the preceding edge. Input zero is the first entry; every other input is
a unique side input. The final edge produces the output, while no source is
an output and no target is an input. All vertices have the same sharing-wire
type. Together with the explicit finite enumerations and source distinctness,
these facts construct the actual finite net.

\lean{Copying/ContractionSpine}{23}{26}{ContractionSpine.mem_sources}
\lean{Copying/ContractionSpine}{28}{30}{ContractionSpine.entry_ne_side}
\lean{Copying/ContractionSpine}{32}{34}{ContractionSpine.entry_injective}
\lean{Copying/ContractionSpine}{36}{39}{ContractionSpine.side_injective}
\par\penalty0
\lean{Copying/ContractionSpine}{41}{43}{ContractionSpine.out_injective}
\par\penalty0
\lean{Copying/ContractionSpine}{45}{54}{ContractionSpine.source_unique}
\lean{Copying/ContractionSpine}{56}{82}{ContractionSpine.source_coverage}
\lean{Copying/ContractionSpine}{84}{98}{ContractionSpine.target_coverage}
\lean{Copying/ContractionSpine}{100}{102}{ContractionSpine.vertex_label}
\par\penalty0
\lean{Copying/ContractionSpine}{105}{125}{ContractionSpine.net}
\par\penalty0

Any adjacency advances strictly in the edge index: a target can meet the
next entry, never a side input or an earlier entry. The generic strict-rank
path theorem therefore gives acyclicity. Every edge is a contraction, so the
finite net satisfies the original contraction-tree definition. Relabelling
its primitives realizes this same tree at any representation allowance.

An attachment map is specified by the successive actual path vertices,
including the selected input and final output, together with one actual side
vertex per contraction. Internal spine vertices map to their corresponding
path vertices; external inputs map to the selected or side attachment; the
output maps to the final path vertex. Direct evaluation proves the entry,
side and output equations.

\lean{Copying/ContractionSpine}{127}{137}{ContractionSpine.adjacent_rank}
\lean{Copying/ContractionSpine}{140}{144}{ContractionSpine.tree}
\lean{Copying/ContractionSpine}{152}{153}{ContractionSpine.atDepth}
\lean{Copying/ContractionSpine}{157}{160}{ContractionSpine.ports}
\par\penalty0
\lean{Copying/ContractionSpine}{162}{175}{ContractionSpine.ports_entry}
\par\penalty0
\lean{Copying/ContractionSpine}{177}{179}{ContractionSpine.ports_side}
\lean{Copying/ContractionSpine}{181}{190}{ContractionSpine.ports_out}

\paragraph{An actual complete-label occurrence.}
Suppose distinct actual contraction edges have exactly those ordered source
rows and single targets. The attachment equations define a
\texttt{MatchEmbedding}: the supplied edge injectivity proves its edge law,
primitive labels agree, and the entry/side/output equations prove both
ordered incidence laws. Because every label is the exact primitive
contraction, the embedding also gives a full recursive-label occurrence.
No injectivity of arbitrary boundary attachments is assumed. Internal
positions and their injectivity follow from the existing complete-incidence
geometry. The selected input and final output coordinates are the supplied
path endpoints by definition.

\lean{Copying/ContractionSpineEmbedding}{20}{36}{ContractionSpine.embedding}
\lean{Copying/ContractionSpineEmbedding}{40}{45}{ContractionSpine.occurrence}
\lean{Copying/ContractionSpineEmbedding}{47}{49}{ContractionSpine.occurrence_input}
\lean{Copying/ContractionSpineEmbedding}{51}{53}{ContractionSpine.occurrence_output}

\paragraph{Extracting the spine from the supplied original tree.}
An exact $n$-step trace has $n+1$ indexed positions. Induction on the trace
constructs this indexing, its first and last equations, each adjacent pair,
and exact prefix and suffix traces. If the successor relation is functional
and the endpoint is terminal, two positions with the same value have suffix
traces of equal length by trace uniqueness. Their indices are therefore
equal. This yields the edge injectivity needed by the spine occurrence.

\lean{Copying/IndexedTrace}{9}{46}{Graph.Trace.indexed}
\lean{Copying/IndexedTrace}{50}{57}{Graph.Trace.indexed_injective}

For an arbitrary original contraction tree, realize it at the requested
depth without changing its incidence graph. A source witness excludes a
weakening and hence identifies a contraction. Every tree edge has exactly
one target, so its target row equals the singleton containing its ordered
target port.

Start with the original selected-input-to-root trace and index its edges.
Use its unique targets as the successive path vertices. The initial input
incidence and subsequent adjacency witnesses show that each path vertex is
consumed by the next indexed edge. Each such edge is a contraction. Its
ordered source row therefore has one actual other source; choose the Boolean
orientation recording which source is selected. The original target rows
and these source rows satisfy the occurrence constructor. The tree's root
output proves that the final coordinate is the original output. The theorem
also retains, for every matched edge, its exact counted original prefix
trace from the selected consuming edge. Thus irrelevant sibling subtrees
are exposed as attachments rather than falsely declared retained, while
all selected edges and their original order are proved to belong to the
supplied tree.

\lean{Copying/MinimalContractionSpine}{14}{15}{ContractionTree.atDepth}
\lean{Copying/MinimalContractionSpine}{18}{26}{ContractionTree.contraction_of_source}
\lean{Copying/MinimalContractionSpine}{29}{36}{ContractionTree.target_port}
\lean{Copying/MinimalContractionSpine}{40}{96}{ContractionTree.minimal_spine}

\subsection{Admissibility of the actual beta cut and copying graft}
\termdef{rooted-complete-cut}
A \emph{rooted complete cut} replaces a fully matched finite source by one
auxiliary hole, retaining every other original edge and incidence. To use
this cut as a specimen context, its binding freedom must be proved. Merely
knowing that the original common context is binding-free is insufficient:
the new hole itself is a new possible endpoint of a forbidden binding path.

First, operation-only reachability excludes a binding ancestor of the reached
wire. Decompose an edge ancestry into its last adjacency. At the reached
wire, target uniqueness identifies the preceding producer with the operation
used by reachability. A zero-length ancestry contradicts that operation's
non-binding label. Otherwise its unique input is the preceding reached wire,
and induction applies. The boundary-input case is impossible because an
input has no producer. This argument permits cycles elsewhere in the graph.
For a reachable term coordinate in a left operational expansion, reflect the
operation path to the common context. A retained operation reflects its
label and ordered incidences. An operation inside a shallow left filling
cannot produce such a reached host output, by the proved one-way obstruction.
Thus actual rooted search supplies original-context reachability, without
adding a new operational premise.
\lean{Copying/RootedBindingBarrier}{9}{16}{Ancestor.last}
\lean{Copying/RootedBindingBarrier}{20}{34}{Search.OperationReachable.no_binding_ancestor}
\lean{Copying/RootedBindingBarrier}{39}{73}{OperationalTemplates.left_reachable_context}

Now suppose every input of a matched source is a reached term wire, every
output is protected by a binding ancestor or has thunk type, and original
holes have their prescribed single term input. Retained adjacencies in the
cut reflect literally to the old graph. An exit from its new hole cannot
lead to an old hole: a protected output supplies an old binding ancestor,
whereas a thunk output must immediately meet a non-hole binding box. Either
case contradicts the original binding freedom. Backwards induction along
an ancestry therefore keeps every path ending at an old hole inside retained
edges.

An edge entering the new hole produces one of the matched inputs. No path
can leave that hole and return to such an input: output protection supplies
a binding ancestor of a reached wire, which the preceding lemma excludes.
A direct self-loop is excluded by the same argument, or by the incompatibility
of a thunk output and a term input. A second backwards induction reduces
paths ending at the new hole to old paths producing an original matched
input. The reachability obstruction excludes every binding source of such
a path. Together these two endpoint cases prove shallow binding freedom of
the cut. Original recursive box bodies are retained unchanged, so their
existing admissibility proves the full recursive property. The beta
application has one reached input; its concrete output guards discharge
these premises. There is no assumed admissibility of an auxiliary frame.
\lean{Copying/RootedCompleteCut}{14}{17}{CompleteMatchCut.cut_adjacent}
\lean{Copying/RootedCompleteCut}{20}{36}{CompleteMatchCut.cut_exit_hole}
\lean{Copying/RootedCompleteCut}{39}{50}{CompleteMatchCut.cut_ancestor_hole}
\lean{Copying/RootedCompleteCut}{53}{57}{CompleteMatchCut.cut_into}
\par\penalty0
\lean{Copying/RootedCompleteCut}{66}{78}{CompleteMatchCut.cut_exit_input}
\lean{Copying/RootedCompleteCut}{83}{95}{CompleteMatchCut.cut_no_loop}
\lean{Copying/RootedCompleteCut}{99}{113}{CompleteMatchCut.cut_ancestor_input}
\lean{Copying/RootedCompleteCut}{117}{151}{CompleteMatchCut.bindingFree_cut}
\par\penalty0

Copying has multiple exposed tree inputs; the argument does not assert that
all of them are rooted. Instead use output guards of the actual original
copy source to prove binding freedom directly after grafting the local
residual context. The earlier guarded replacement argument needed a wholly
hole-free replacement. Here only \emph{shallow} holes must be absent. Any
forbidden path to a surviving outer hole reflects to the original graph,
where the guard gives a contradiction. Nested holes inside replacement boxes
are allowed and checked by the separately proved recursive binding freedom
of the replacement. This is exactly the distinction required for copying
payloads with arbitrary deep template occurrences.
\lean{Copying/BindingOpenGuardedReplacement}{15}{35}{CompleteMatchCut.guarded_shallow_binding_open}
\lean{Copying/BindingOpenGuardedReplacement}{39}{55}{CompleteMatchCut.bindingFree_replace_open_guarded}

\subsection{Extracting the local deep transformation from an actual match}
The following construction starts with an actual complete occurrence
$f:A\to C[L]$ of a local rule source. Every outer edge of this occurrence
has already been proved to avoid a shallow operational hole of the original
common context $C$. Its contribution may therefore be retained, or may be
an original box whose body contains an arbitrary recursive filling. The
construction below derives the local decomposition needed by the proof from
these certificates. No decomposition of the occurrence into a preselected
frame is assumed.

\paragraph{Retained or recursively boxed contributions.}
A retained atomic certificate has two alternatives. The first supplies an
actual recursive isomorphism between its endpoints. The second supplies an
isomorphism from the left endpoint to a boxed body, a genuine filling
witness between the two bodies, binding freedom of that witness's actual
common context, and an isomorphism from the other box to the right endpoint.
Its interface equalities preserve the box's bound-variable and free-variable
interfaces. These are finite structural certificates, with no observational
equivalence included in their definition.

The auxiliary exact-context witness constructor converts an original
recursive filling into a filling witness while retaining a proof that the
witness's stored context is exactly the supplied original context. This
equality is needed to reuse its binding-freedom proof. The general witness
constructor, which selects an existential witness, alone would not identify
that stored context. Endpoint congruence composes the existing isomorphisms;
interface transport changes only the dependent indices.

\lean{Copying/RetainedAtomic}{12}{17}{FillingWitness.ofDeepFillingExact}

\lean{Copying/RetainedAtomic}{20}{29}{RetainedAtomic}

\lean{Copying/RetainedAtomic}{34}{38}{RetainedAtomic.congr}

\lean{Copying/RetainedAtomic}{40}{43}{RetainedAtomic.interfaces}

A primitive footprint cannot have the recursively boxed alternative, because
its outer edge labels are primitive. Hence its retained atomic certificate
is an isomorphism. For a boxed footprint, the box view exposes either that
isomorphism or the actual paired body and its binding-free common context.
The view is obtained by inspecting the two certificate constructors and the
recursive label equivalence. Finally, an original edge-filling certificate
at a non-hole edge supplies a retained atomic certificate: inspect its
retained or boxed constructor, rule out the filled-hole constructor, and
use the original context's recursive binding freedom for the boxed body.

\lean{Copying/RetainedAtomic}{45}{54}{RetainedAtomic.primitiveIso}

\lean{Copying/RetainedAtomic}{56}{60}{RetainedAtomic.BoxView}

\lean{Copying/RetainedAtomic}{62}{77}{RetainedAtomic.boxView}

\lean{Copying/RetainedAtomic}{81}{97}{RetainedAtomic.of_edge}

\paragraph{Decoration of the actual local source.}
A retained decoration assigns to every edge of $A$ a right replacement and
a retained atomic certificate from that edge's original corolla to the
replacement. Its graph is the simultaneous expansion of $A$ by this family.
If $A$ is entirely primitive, the atomic isomorphisms assemble into an
isomorphism from $A$ to this expansion; its vertex equation preserves each
original host coordinate exactly. Pulling a decoration along a recursive
isomorphism reindexes its edges and interfaces. The resulting expansion
is isomorphic to the original decorated graph, with the stated equation on
all host ports.

\lean{Copying/RetainedDecoration}{10}{13}{RetainedDecoration}

\lean{Copying/RetainedDecoration}{17}{17}{RetainedDecoration.graph}

\lean{Copying/RetainedDecoration}{19}{22}{RetainedDecoration.primitiveIso}

\lean{Copying/RetainedDecoration}{24}{26}{RetainedDecoration.primitiveIso_vertex}

\par\penalty0
\lean{Copying/RetainedDecoration}{28}{35}{RetainedDecoration.pull}

\lean{Copying/RetainedDecoration}{37}{45}{RetainedDecoration.pullIso}

\lean{Copying/RetainedDecoration}{47}{49}{RetainedDecoration.pullIso_host}

A decoration of an expanded graph restricts to each original component.
Reassociating these simultaneous expansions gives the grouped isomorphism;
for a single corolla it reduces to the unique replacement graph. Serial and
tensor restrictions select their original left and right components. Their
assembly isomorphisms follow from expansion associativity and the ordered
parallel-construction isomorphism. These operations inspect the existing
local source graph and preserve its ordered interfaces.

\lean{Copying/RetainedDecoration}{51}{54}{RetainedDecoration.component}

\lean{Copying/RetainedDecoration}{56}{58}{RetainedDecoration.groupedIso}

\lean{Copying/RetainedDecoration}{60}{63}{RetainedDecoration.corollaIso}

\lean{Copying/RetainedDecoration}{65}{67}{RetainedDecoration.serialLeft}

\par\penalty0
\lean{Copying/RetainedDecoration}{69}{71}{RetainedDecoration.serialRight}

\lean{Copying/RetainedDecoration}{73}{78}{RetainedDecoration.serialIso}

\lean{Copying/RetainedDecoration}{80}{83}{RetainedDecoration.tensorLeft}

\lean{Copying/RetainedDecoration}{85}{88}{RetainedDecoration.tensorRight}

\par\penalty0
\lean{Copying/RetainedDecoration}{90}{97}{RetainedDecoration.tensorIso}

The actual occurrence now determines a decoration. For a local edge $e$,
use the original common-context owner of $f(e)$, choose its certified right
contribution, and transport its interface to the interface of $e$. The
non-hole premise supplies the retained atomic certificate. The skeleton
map has identity local carriers and changes each label to the original
owner label. Simultaneous-expansion transport identifies the decorated
local graph with the right expansion of this very skeleton. Its host-port
equation is exact, so a local right computation will occur at the required
original context coordinate.

\lean{Copying/RetainedDecoration}{108}{114}{RetainedFootprint.decoration}

\lean{Copying/RetainedDecoration}{116}{122}{RetainedFootprint.decorationSkeleton}

\lean{Copying/RetainedDecoration}{126}{135}{RetainedFootprint.decorationIso}

\lean{Copying/RetainedDecoration}{137}{141}{RetainedFootprint.decorationIso_host}

\paragraph{Lambda views and shallow-hole control.}
Complete recursive-label occurrences reflect any global primitive-label
property, by reflecting it edge by edge along label equivalence. In
particular, the actual local source inherits focus freedom and hole freedom
from the original expanded graph. A lambda view specializes the boxed view
to the original translated lambda box. The equality $k+1=1+k$ is used only
to align its indexed interface; the body and the term translation are
unchanged, and the existing common-context certificate is transported with
its binding-freedom proof.

\lean{Copying/RetainedLocalViews}{11}{13}{Opening.Occurrence.allLabels}

\lean{Copying/RetainedLocalViews}{19}{23}{RetainedAtomic.LambdaView}

\lean{Copying/RetainedLocalViews}{25}{42}{RetainedAtomic.lambdaView}

Absence of shallow holes is weaker than recursive hole freedom. Serial
composition preserves the former by inspecting the component owning each
edge. Recursive hole freedom implies it directly. The copy-construction
lemmas apply this fact to the active and retained copies, the residual
sharing tree, and the unchanged output distributors. In the deep-copy
witness, all changed body holes are still inside original boxes. Thus its
whole local common context has no outer hole, even though its recursive
box bodies may contain arbitrary operational holes. This precise property
will establish admissibility of the global copying residual.

\lean{Copying/RetainedLocalViews}{48}{55}{FillingWitness.serial_no_shallow_holes}

\lean{Copying/RetainedLocalViews}{57}{62}{FillingWitness.holeFree_no_shallow}

\lean{Copying/RetainedLocalViews}{64}{79}{FillingWitness.copies_no_shallow_holes}

\lean{Copying/RetainedLocalViews}{81}{87}{FillingWitness.copyingAfter_no_shallow_holes}

\par\penalty0
\lean{Copying/RetainedLocalViews}{89}{93}{FillingWitness.deepCopy_no_shallow_holes}

\subsection{Constructed local beta and copying residuals}
\paragraph{Beta.}
Restrict the decoration of the original beta source to its lambda box.
All remaining components are primitive, including the original stable
argument. Their decoration is therefore isomorphic to the original
components. Serial and tensor assembly give an isomorphism from the actual
paired local source to the unchanged application shape with only its box
component varied. Inspect the lambda view. In the retained case, transport
the original beta step along this isomorphism and use a reflexive witness
on its reduct. In the deep case, use the original beta step with the paired
body and the same stable argument, then use the previously proved deep-beta
body residual. Its common context is binding-free. Both alternatives return
an actual machine rewrite, resuming at the original sole input.

\lean{Copying/OperationalDeepLocal}{14}{14}{OperationalTemplates.DeepLocal.Beta.boxes}

\lean{Copying/OperationalDeepLocal}{16}{24}{OperationalTemplates.DeepLocal.Beta.sourceIso}

\lean{Copying/OperationalDeepLocal}{28}{45}{OperationalTemplates.DeepLocal.Beta.transport}

\paragraph{Copying.}
For a primitive copyable payload the complete decorated source is
isomorphic to the original source, so the original copying step transports
directly and its local residual is reflexive. Focus and hole preservation
make that residual binding-free and free of shallow holes. For a lambda
payload, restrict to its box and reconstruct the unchanged tree, lambda
edge and box interface. The retained box case is again reflexive. The deep
case copies the actual paired body witness and uses the previously proved
binding-free deep-copy construction. Its absence of shallow holes follows
from the lemmas above.

Induction over the original copyability proof covers instances, constants,
application, unary and binary operation payloads, lambdas, and recursive
representation changes. The last case pulls the decoration through the
original representation isomorphism, invokes the induction hypothesis,
and transports the actual source and reduct back. Thus active operation
payloads and all original copyable constructors remain included.

\lean{Copying/OperationalDeepLocal}{51}{62}{OperationalTemplates.DeepLocal.Copying.primitive_transport}

\lean{Copying/OperationalDeepLocal}{68}{68}{OperationalTemplates.DeepLocal.Copying.boxes}

\lean{Copying/OperationalDeepLocal}{70}{76}{OperationalTemplates.DeepLocal.Copying.sourceIso}

\lean{Copying/OperationalDeepLocal}{78}{97}{OperationalTemplates.DeepLocal.Copying.lambda_transport}

\par\penalty0
\lean{Copying/OperationalDeepLocal}{101}{136}{OperationalTemplates.DeepLocal.Copying.transport}

\paragraph{Extraction and admissibility.}
The open guarded graft theorem applies to the actual common skeleton cut.
Transport the original matched source's binding-output guards to its
skeleton, then insert the local witness context. Its binding freedom and
absence of shallow holes permit arbitrary recursive holes inside boxes.
The original host's binding freedom excludes every path from a sourced
binding edge to a surviving shallow hole; the output guards handle paths
that leave the inserted component. No binding-freedom premise is imposed
on the intermediate cut in this argument.

The beta and copying extraction theorems apply the local constructions to
the decoration obtained from the actual occurrence. Focus and hole freedom
are reflected from the expanded host. Transport along the decoration-to-
skeleton isomorphism then gives a right local rewrite at the exact
normalized host port. Beta returns a binding-free local witness. Copying
also returns its proved absence of shallow holes. The source match is
arbitrary subject to the already established non-hole footprint condition.

\lean{Copying/OperationalDeepExtraction}{16}{25}{RetainedFootprint.residualGraft_binding_open_guarded}

\lean{Copying/OperationalDeepExtraction}{36}{50}{OperationalTemplates.DeepExtraction.beta}

\lean{Copying/OperationalDeepExtraction}{56}{70}{OperationalTemplates.DeepExtraction.copying}

\subsection{Whole-graph diamonds and full forward robustness}
\paragraph{Exact residual positions.}
Suppose the left and right local steps resume at the same numbered input,
and a genuine paired local witness has a binding-free geometric graft.
Complete-cut transfer constructs both whole-graph rewrites. The pointed
residual construction inserts that local common input in the actual graft.
A filling witness fixes its input coordinates on both sides, so this point
is exactly the two computed resume ports. Rewrite preservation gives
closed endpoint graphs. The resulting pair is an original binding-free
quasi-specimen with equality on both sides.

\lean{Copying/OperationalCopyResidual}{14}{38}{RetainedFootprint.input_diamond}

For copying, the exhaustive forward overlap theorem supplies footprint
avoidance. Extract the actual local paired rewrite and witness, use the
original copyable source's output guards and the open guarded graft theorem,
and apply the input diamond. This yields one actual whole-graph rewrite on
each side and the required original quasi-specimen, for every finite sharing
tree and every copyable payload. The quantitative theorem compares an
independently supplied left successor with the constructed one using the
proved machine determinism, takes zero additional left steps and one right
step, and weakens endpoint equalities to the requested refinements.

\lean{Copying/OperationalCopyResidual}{51}{66}{OperationalTemplates.ForwardResidual.copying_diamond}

\lean{Copying/OperationalCopyResidual}{70}{84}{OperationalTemplates.ForwardResidual.copying}

\paragraph{Why beta needs its rooted cut.}
Opening a box can expose shallow holes in the body's common context, so
beta uses a different admissibility argument. A binding-free local witness
can be plugged into a binding-free cut at a sourced hole; interface
reindexing preserves its context certificate. To establish binding freedom
of the actual beta cut, use the original left specimen's search-rootedness.
Search reachability reflects to the original common context, and a reachable
term position has no binding-edge ancestor. The source's sole input is
therefore protected. Its free closure outputs have the original box as
binding ancestor; the unchanged stable argument supplies either a guarded
output or a thunk-typed output. The proved rooted-cut theorem applies to
these precise ports. Thus cut admissibility is derived from the original
rootedness premise, not added to the semantics.

The beta diamond combines this cut with the extracted deep local residual
and the input diamond. Its quantitative corollary has the same zero-left,
one-right continuation counts as copying. The full original body, stable
argument and ordered output interfaces occur in these statements.

\lean{Copying/OperationalBetaResidual}{15}{26}{RetainedFootprint.residualGraft_binding_frame}

\lean{Copying/OperationalBetaResidual}{41}{61}{OperationalTemplates.ForwardResidual.beta_frame_binding}

\lean{Copying/OperationalBetaResidual}{65}{80}{OperationalTemplates.ForwardResidual.beta_diamond}

\lean{Copying/OperationalBetaResidual}{84}{98}{OperationalTemplates.ForwardResidual.beta}

\paragraph{Universal forward assembly.}
\termdef{operational-forward-robustness}
\emph{Operational forward robustness} is the original specimen-system
robustness property for either micro-beta or reference creation, the
original binding-free context class, and the full copying UAM. It quantifies
every rooted remote rewrite specimen and every successor of its left state.
The exhaustive local-case theorem inspects all nine original local rule
constructors: allocation, read, assignment, the two atom-equality rules,
unary and binary arithmetic, beta and copying. Each branch invokes its
constructed whole-graph residual result. The operational stability and
copyability premises are used exactly as supplied by the original rules.

An arbitrary original contextual rewrite contains a complete local
occurrence, so the rewrite-case theorem requires no occurrence or footprint
certificate from its caller. Reflection of a state-class step to the
specified finite representative then supplies such a rewrite; search has
no transition from rewrite mode. Finally, unfold the original specimen
coordinates, extract their actual rootedness, binding freedom, closed
fillings and non-entering token, and apply this exhaustive dispatch.
The forward robustness theorem follows for arbitrary cost orders $Q,P,U$.
Each constructed residual case uses $k=0$ and $\ell=1$, so its cost
comparison uses the reflexivity of $Q$ at $1$.

\lean{Copying/FullOperationalForwardRobustness}{22}{42}{OperationalTemplates.ForwardResidual.local_case}

\lean{Copying/FullOperationalForwardRobustness}{46}{57}{OperationalTemplates.ForwardResidual.rewrite_case}

\lean{Copying/FullOperationalForwardRobustness}{61}{74}{OperationalTemplates.ForwardResidual.class_step}

\lean{Copying/FullOperationalForwardRobustness}{79}{108}{OperationalTemplates.template_robust}

\paragraph{Precisely what the reverse residual results establish.}
The copying and beta diamonds also yield reverse quantitative matching for
every independently represented right successor: exchange the two actual
steps and reverse the original quasi-specimen. These two theorems retain an
actual original \emph{left} occurrence; beta retains the same original left
rootedness used to justify its cut. They do not themselves extract a left
match from an arbitrary right rewrite. Consequently they must not be read
as the universal converse robustness theorem, or as Proposition~12.2.

\lean{Copying/OperationalDeepConverse}{18}{33}{OperationalTemplates.ForwardResidual.copying_converse}

\lean{Copying/OperationalDeepConverse}{48}{63}{OperationalTemplates.ForwardResidual.beta_converse}

\subsection{Extending and extracting minimal copying occurrences}
A complete occurrence $f:A\to B$ with the same output interface extends
through serial composition with an unchanged graph $H$, provided that $f$
fixes each output port. Its input ports need not remain external: they may
attach internally in $B$. The extended port map keeps the common serial
interface and the internals of $H$, maps the internals of $A$ through $f$,
and sends external input attachments through the original left serial
insertion. Its edge map applies $f$ to left edges and keeps right edges.
The two insertion equations follow by cases on internal, input and output
ports; the shared-output hypothesis is used only on left outputs.

\lean{Copying/SerialOccurrenceExtension}{13}{20}{Opening.Occurrence.SerialExtension.ports}

\lean{Copying/SerialOccurrenceExtension}{22}{26}{Opening.Occurrence.SerialExtension.edges}

\lean{Copying/SerialOccurrenceExtension}{31}{48}{Opening.Occurrence.SerialExtension.left_insert}

\lean{Copying/SerialOccurrenceExtension}{51}{68}{Opening.Occurrence.SerialExtension.right_insert}

The extended edge map is injective because the serial components have
disjoint edge carriers and $f$ is injective. The insertion equations
transport every ordered source and target list. Together with recursive
label preservation, these facts produce a complete occurrence and its
induced injective internal map. The left and right coordinate theorems
record the two original insertions exactly. If one selected input of $f$
maps to a specified external input, the extension preserves that numbered
external input as well.

\lean{Copying/SerialOccurrenceExtension}{70}{118}{Opening.Occurrence.SerialExtension.matchEmbedding}

\lean{Copying/SerialOccurrenceExtension}{122}{129}{Opening.Occurrence.SerialExtension.occurrence}

\lean{Copying/SerialOccurrenceExtension}{131}{133}{Opening.Occurrence.SerialExtension.occurrence_left}

\lean{Copying/SerialOccurrenceExtension}{135}{137}{Opening.Occurrence.SerialExtension.occurrence_right}

\par\penalty0
\lean{Copying/SerialOccurrenceExtension}{142}{146}{Opening.Occurrence.SerialExtension.occurrence_input}

Apply this extension to the minimal oriented spine extracted from an
arbitrary original contraction tree. Its selected input is the spine's
first input and its sole output is the original tree output. Serially
compose with the entire original payload $H$. The result is a complete
occurrence of a smaller original copying source inside the supplied
source. The selected input is preserved, every payload port is unchanged,
and every retained spine contraction has its original counted ancestry
from the active input. Unused sibling branches remain attachment ports of
the smaller occurrence. This is a proved source-rematching construction;
its statement does not assume or assert a new compute rule for truncated
trees, and does not alter the original copying semantics.

\lean{Copying/MinimalCopyOccurrence}{14}{40}{CopyRule.minimal_occurrence}

\subsection{Converse classification from rooted history}
\termdef{converse-selected-match}
A \emph{converse selected match} records the actual information inspected by a
rewrite on the right of an original operational specimen, together with its
origin in the common context. For allocation, arithmetic and beta, the whole
local occurrence avoids shallow template holes. For store access, the selected
operation, atom trace and store are reflected. For copying, the selected
sharing trace and complete payload are reflected. The last two cases retain
certificates for the supplied whole successor, including its resumed token.
This classification covers every original contextual rewrite. Turning the
reflected store and copying data into matching original computations and
admissible residual specimens is a further obligation for universal converse
robustness.

Fix a common context $C$, original paired recursive fillings $L,R$, and a
common vertex $v$. Write $\iota_L,\iota_R$ for the two host-coordinate maps.
The token has rewrite mode at $v$ and is not entering a shallow hole. Assume
that the original left state at $\iota_L(v)$ has a rooted search history.
Binding freedom of $C$ is required for the store and copying path arguments.
The declaration names below are in \texttt{OperationalTemplates.ConverseOverlap}
unless another namespace is displayed.

\paragraph{Rooted argument readiness.}
For a fixed graph $G$, $\operatorname{StableAt}(G,u)$ means that some actual
edge consumes $u$ and carries a stable primitive label: a passive operation
or an instance edge. It asserts a local source property, rather than stability
of the whole graph.
\lean{Copying/ConverseArgumentHistory}{11}{12}{Search.StableAt}

The inductive readiness stack starts at a graph input. An argument frame
records an operation with its singleton source, its parent frame, the selected
ordered argument $j$, and stability of every argument preceding $j$. This is
an invariant used to prove a property of the existing search relation; it
adds no stack to the machine state. The parent lemma recovers the precise
frame from the current argument. Target incidence uniqueness identifies the
operation edge, its label identifies the operation, singleton source equality
identifies the parent wire, and duplicate-free target lists identify $j$.
\lean{Copying/ConverseArgumentHistory}{16}{22}{Search.ReadyStack}
\lean{Copying/ConverseArgumentHistory}{26}{43}{Search.ReadyStack.parent}

The mode-dependent invariant combines this stack with two assertions:
backtracking has a stable source at the current wire, and a rewrite request
at an operation has a stable source at every evaluated argument position.
\lean{Copying/ConverseArgumentHistory}{46}{50}{Search.ReadyHistory}

Preservation is proved by cases on the original search step. Entering an
operation with arguments pushes its first argument, whose preceding prefix
is empty. Passive completion and instance entry supply a stable source.
Backtracking recovers the parent frame and adds the just-completed argument
to the stable prefix. Search then visits the next argument or returns to the
operation with every argument ready. A contraction rewrite cannot also be an
operation rewrite at the same wire, by source uniqueness and incompatible
labels. Induction on a finite search run initializes this invariant at the
external input. Thus every operation rewrite reached by an actual rooted
search has all its arguments ready in that same unchanged graph.
\lean{Copying/ConverseArgumentHistory}{53}{130}{Search.readyHistory_step}
\lean{Copying/ConverseArgumentHistory}{133}{142}{Search.readyHistory_run}
\lean{Copying/ConverseArgumentHistory}{146}{151}{Search.rooted_rewrite_arguments}

\paragraph{Reversing certificates without changing their context.}
The converse pre-template has relation $T^{\mathrm{op}}(G,H)=T(H,G)$;
its focus-freedom proof exchanges the two original endpoint proofs.
Induction on the box allowance reverses a recursive filling while preserving
its literal common context. Retained label copies are exchanged, filled
holes use the reversed relation, and boxed components recursively exchange
their body fillings. The corresponding edge theorem allows the generic
binding-path barriers to inspect the original right expansion. No symmetry
assumption is imposed on the directed operational template.
\lean{Copying/ConverseFillings}{8}{10}{PreTemplate.converse}
\lean{Copying/ConverseFillings}{13}{32}{DeepFilling.converse}
\lean{Copying/ConverseFillings}{36}{41}{DeepFillingEdge.converse}

\paragraph{Excluding shallow right contributions at evaluated coordinates.}
Suppose the left expansion has a stable source at $\iota_L(u)$ and a right
edge consumes $\iota_R(u)$. Both sources have the same original owner by
source incidence uniqueness in $C$. If that owner were an operational hole,
the corresponding left input would expose the active application or reference
head. It cannot also expose a stable primitive. The earlier left-head
exclusion handles retained unfilled holes as well.
\lean{Copying/ConverseHeadOrigin}{16}{26}{OperationalTemplates.ConverseOverlap.stable_source_not_hole}

At the non-entering token, the actual right operation head is retained. Its primitive
label and singleton source therefore also occur on the left, with the same
original owner. Retained target rows identify each ordered right argument as
$\iota_R(u)$ for an original $u$ and identify its left counterpart as
$\iota_L(u)$. Applying rooted argument readiness to the left head proves
stability at that counterpart. Consequently any actual right edge consuming
the selected argument avoids every shallow operational hole, independently
of its own full label.
\lean{Copying/ConverseHeadOrigin}{30}{43}{OperationalTemplates.ConverseOverlap.left_head}
\lean{Copying/ConverseHeadOrigin}{47}{65}{OperationalTemplates.ConverseOverlap.argument_stable}
\lean{Copying/ConverseHeadOrigin}{69}{77}{OperationalTemplates.ConverseOverlap.argument_source_not_hole}

There is a separate typed-interface exclusion. Both operational families have
exactly one term input. An inspected non-hole edge consuming a non-term wire
at an original coordinate therefore cannot come from a shallow operational
filling. A retained hole is excluded by the inspected edge's own non-hole
label. This covers the thunk-consuming box reached after a lambda. Finally,
every output of an active operation in the example alphabet is an evaluated
term argument, so adjacency after an active head is a convenient specialization
of argument exclusion.
\lean{Copying/ConverseHeadOrigin}{81}{101}{OperationalTemplates.ConverseOverlap.nonterm_source_not_hole}
\lean{Copying/ConverseHeadOrigin}{104}{116}{OperationalTemplates.ConverseOverlap.adjacent_source_not_hole}

\paragraph{Complete allocation, arithmetic and beta footprints.}
The elementary footprint lemma treats any source whose edges are its active
head or direct successors of that head. The non-entering condition excludes
the head from shallow holes; rooted argument exclusion handles every successor.
This proves a conclusion about the entire supplied occurrence.
\lean{Copying/ConversePrimitiveOverlap}{11}{26}{OperationalTemplates.ConverseOverlap.simple_avoids}

Inspecting the allocation source shows that it contains the reference head
and the original stable operand. Each edge of that operand consumes its sole
input, so it is adjacent to the reference head. The resulting exclusion
preserves the stable operand and its complete interface, including thunk
outputs. Unary arithmetic contains its head and one numeric constant;
binary arithmetic contains its head and the two constants at their original
ordered positions. Their concrete edge classifications instantiate the same
lemma for the original arithmetic rules.
\lean{Copying/ConversePrimitiveOverlap}{29}{40}{OperationalTemplates.ConverseOverlap.allocation_edges}
\lean{Copying/ConversePrimitiveOverlap}{43}{52}{OperationalTemplates.ConverseOverlap.allocation_avoids}
\lean{Copying/ConversePrimitiveOverlap}{55}{65}{OperationalTemplates.ConverseOverlap.unary_edges}
\lean{Copying/ConversePrimitiveOverlap}{68}{82}{OperationalTemplates.ConverseOverlap.binary_edges}
\par\penalty0
\lean{Copying/ConversePrimitiveOverlap}{85}{93}{OperationalTemplates.ConverseOverlap.unary_avoids}
\lean{Copying/ConversePrimitiveOverlap}{96}{105}{OperationalTemplates.ConverseOverlap.binary_avoids}

For beta, the application head is retained by non-entry. Its function wire
leads to a lambda whose source is an evaluated argument, hence whose original
owner is not a hole. The lambda's retained target supplies the original
coordinate of the box input. That wire has thunk type, so the typed-interface
exclusion retains the box as well. Every edge of the stable argument consumes
the application's other evaluated argument. These cases exhaust the outer
beta footprint. The proof keeps the actual recursive box label and does not
inspect or restrict the lambda body.
\lean{Copying/ConverseBetaOverlap}{13}{72}{OperationalTemplates.ConverseOverlap.beta_avoids}

\paragraph{Selected store traces and intrinsic matches.}
For a retained edge, its entire target list is the original target list mapped
by $\iota_R$. This identifies actual stored-content and value coordinates.
A retained atom hop reflects both its forward label and its adjacency into
$C$. Induction on an atom trace then reflects every hop after a sourced
binding ancestor: binding freedom prevents the next owner from being a
shallow hole. The trace length is preserved exactly. These lengths count
atom-resolution hops, not additional machine transitions.
\lean{Copying/ConverseStoreOverlap}{16}{19}{OperationalTemplates.ConverseOverlap.retained_target}
\lean{Copying/ConverseStoreOverlap}{23}{34}{OperationalTemplates.ConverseOverlap.store_hop_owner}
\lean{Copying/ConverseStoreOverlap}{38}{50}{OperationalTemplates.ConverseOverlap.store_trace_owner}

A sourced initial atom hop is binding. Its forward label is an instance or
atom contraction, or it would be a weakening; the last possibility contradicts
the nonempty source witness because weakening has no inputs. To reflect an
operation operand, first retain its initial edge by rooted argument exclusion.
A zero-hop resolution is already retained. For a nonempty resolution, use the
binding initial edge and the preceding trace induction. Reflect the final
store label and the original ordered argument lookup. This constructs the
same operand resolution in $C$, with the same store owner and hop count, and
also proves that the store owner is not a hole. No acyclicity of the surrounding
store graph is assumed.
\lean{Copying/ConverseStoreOverlap}{54}{69}{OperationalTemplates.ConverseOverlap.store_start_binding}
\lean{Copying/ConverseStoreOverlap}{73}{106}{OperationalTemplates.ConverseOverlap.operand_owner}

Two actual retained edges have equal original owners exactly when they are
the same edge: each retained component has only its specified corolla edge.
Reflecting both operands therefore preserves both equality and disequality
of their stores, giving the identical Boolean reference-identity decision in
$C$.
\lean{Copying/ConverseStoreOverlap}{110}{119}{OperationalTemplates.ConverseOverlap.owner_eq_iff}
\lean{Copying/ConverseStoreOverlap}{123}{131}{OperationalTemplates.ConverseOverlap.equality_owner}

A read match is reconstructed from its retained head, reflected operand,
retained store and singleton store-target row. It preserves the actual store
identity, lookup count and content coordinate. A write match additionally
reflects the head's ordered instance and new-value wires and the store's old
content wire. Injectivity of the host-coordinate map recovers the exact target
lists. These are intrinsic matches determined by actual incidence; no atom-tree
presentation is supplied as a new assumption.
\lean{Copying/ConverseStoreOverlap}{135}{158}{OperationalTemplates.ConverseOverlap.read_match_owner}
\lean{Copying/ConverseStoreOverlap}{162}{198}{OperationalTemplates.ConverseOverlap.assignment_match_owner}

\paragraph{Exhaustive classification and its remaining obligation.}
The inductive classification has five outcomes: a complete local occurrence
avoiding shallow holes; a reflected read match; a reflected write match; a
reflected reference-identity decision; or a copying certificate with retained
payload and reflected resolution. The complete-occurrence outcome also identifies the supplied whole successor
with the explicit finite-cut replacement, including its resumed token. The
read and write outcomes record exact coordinate and count correspondences. Their canonical successor is recursively
isomorphic to the actual supplied successor, including the resumed token.
The equality outcome records its Boolean decision and the same successor
identification. The copy outcome retains the original whole-graph realization
certificate, the entire normalized payload occurrence and the exact selected
resolution in $C$.
\lean{Copying/ConverseClassification}{14}{49}{OperationalTemplates.ConverseOverlap.Case}

For reads, writes and equality, extract the intrinsic source information from
the original contextual rewrite, reflect it by the preceding theorems, and
apply the already proved canonical-successor correspondence. For copying,
extract the original complete successor certificate. Apply the binding-path
barrier to the reversed filling certificate: every outer payload edge,
including a lambda box, lies beyond the retained active contraction, and the
selected sharing trace reflects with its exact length. This copying argument
does not require the left rooted-history premise.
\lean{Copying/ConverseClassification}{60}{64}{OperationalTemplates.ConverseOverlap.classify_read}
\lean{Copying/ConverseClassification}{67}{71}{OperationalTemplates.ConverseOverlap.classify_assignment}
\lean{Copying/ConverseClassification}{74}{78}{OperationalTemplates.ConverseOverlap.classify_equality}
\lean{Copying/ConverseClassification}{83}{91}{OperationalTemplates.ConverseOverlap.classify_copy}
\par\penalty0

Finally, extract a full-label local occurrence from an arbitrary actual right
rewrite, including any sequence of original renaming, plugging and expansion
constructors. Case analysis covers all nine original local constructors:
allocation, read, assignment, same-store equality, different-store equality,
unary arithmetic, binary arithmetic, beta and copying. Allocation, the two
arithmetic families and beta use the complete-footprint outcome; the remaining
constructors use their corresponding intrinsic outcomes. Thus classification
requires no overlap premise and no caller-provided residual certificate.
\lean{Copying/ConverseClassification}{95}{116}{OperationalTemplates.ConverseOverlap.rewrite_classification}

The theorem deliberately makes no claim that every unused sibling of an
arbitrarily large supplied contraction tree has a non-hole owner. Its store
and copying outcomes concern the selected traces and data actually used by
the rule. Completing universal converse robustness still requires realizing
suitable original local matches from these data, composing their whole-graph
responses and proving binding freedom of the residual specimens. In particular,
the intrinsic selector certificate alone is not a proof that the common
context performs an original machine transition. The terms of Proposition~12.2,
their translation, the operation alphabet, and the stability and copyability
premises remain unchanged.
